\element{EM1.inv}{
  \begin{questionmult}{invarianceBoule}\bareme{haut=2,p=0}
    En coordonnées sphériques, le champ électrique créé par une boule uniformément chargée à l'origine du repère dépend de
    \begin{multicols}{4}
      \begin{reponses}
        \bonne{$r$}
        \mauvaise{$\theta$}
        \mauvaise{$\phi$}
        \mauvaise{$x$}
        \mauvaise{$y$}
        \mauvaise{$z$}
      \end{reponses}
    \end{multicols}
  \end{questionmult}
}
\element{EM1.inv}{
  \begin{questionmult}{invarianceCylindre}\bareme{haut=2,p=0}
    En coordonnées cylindriques, le champ électrique créé par un cylindre infini uniformément chargé dont l'axe est confondu avec $(Oz)$ dépend de
    \begin{multicols}{4}
      \begin{reponses}
        \bonne{$r$}
        \mauvaise{$\theta$}
        \mauvaise{$\phi$}
        \mauvaise{$x$}
        \mauvaise{$y$}
        \mauvaise{$z$}
      \end{reponses}
    \end{multicols}
  \end{questionmult}
}
\setgroupmode{EM1.inv}{cyclic}




\element{EM1.sym}{
  \begin{question}{symetrieBoule}\bareme{1}
    En coordonnées sphériques, le champ électrique créé par une boule uniformément chargée à l'origine du repère est dirigé selon
    \begin{multicols}{6}
      \begin{reponses}
        \bonne{$\vv{e_r}$}
        \mauvaise{$\vv{e_\theta}$}
        \mauvaise{$\vv{e_\phi}$}
        \mauvaise{$\vv{e_x}$}
        \mauvaise{$\vv{e_y}$}
        \mauvaise{$\vv{e_z}$}
      \end{reponses}
    \end{multicols}
  \end{question}
}
\element{EM1.sym}{
  \begin{question}{symetrieCylindre}\bareme{1}
    En coordonnées cylindriques, le champ électrique créé par un cylindre infini uniformément chargé dont l'axe est confondu avec $(Oz)$ est dirigé selon
    \begin{multicols}{6}
      \begin{reponses}
        \bonne{$\vv{e_r}$}
        \mauvaise{$\vv{e_\theta}$}
        \mauvaise{$\vv{e_\phi}$}
        \mauvaise{$\vv{e_x}$}
        \mauvaise{$\vv{e_y}$}
        \mauvaise{$\vv{e_z}$}
      \end{reponses}
    \end{multicols}
  \end{question}
}
\setgroupmode{EM1.sym}{cyclic}



\element{EM1.surfGauss}{
  \begin{question}{surfaceGaussCylindre}\bareme{1}
    En coordonnées cylindriques, le champ créé par un cylindre uniformément chargée s'écrit $\vv{E}=E(r)\vv{e_r}$. La surface de Gauss adaptée à l'étude est
    \begin{multicols}{3}
      \begin{reponses}
        \mauvaise{Un disque}
        \mauvaise{Un cercle}
        \mauvaise{Une sphère}
        \mauvaise{Un cube}
        \bonne{Un cylindre fermé (avec ses deux extrémités)}
        \mauvaise{Un cylindre ouvert (sans ses deux extrémités)}
        \mauvaise{Une boule}
      \end{reponses}
    \end{multicols}
  \end{question}
}
\element{EM1.surfGauss}{
  \begin{question}{surfaceGaussBoule}\bareme{1}
    En coordonnées sphériques, le champ créé par une boule uniformément chargée s'écrit $\vv{E}=E(r)\vv{e_r}$. La surface de Gauss adaptée à l'étude est
    \begin{multicols}{3}
      \begin{reponses}
        \mauvaise{Un disque}
        \mauvaise{Un cercle}
        \bonne{Une sphère}
        \mauvaise{Un cube}
        \mauvaise{Un cylindre fermé (avec ses deux extrémités)}
        \mauvaise{Un cylindre ouvert (sans ses deux extrémités)}
        \mauvaise{Une boule}
      \end{reponses}
    \end{multicols}
  \end{question}
}
\setgroupmode{EM1.surfGauss}{cyclic}




\element{EM1.Quint}{
  \begin{question}{QintBouleInt}\bareme{1}
    La charge contenue dans une surface de Gauss sphérique de rayon $r\leq R$ entourant une boule de rayon $R$ uniformément chargé de densité volumique de charge $\rho$ est
    \begin{multicols}{3}
      \begin{reponses}
        \bonne{$\frac{4}{3}\pi r^3$}
        \mauvaise{$\frac{4}{3}\pi R^3$}
        \mauvaise{$4\pi r^2$}
        \mauvaise{$4\pi R^2$}
        \mauvaise{$4\pi r^3$}
        \mauvaise{$4\pi R^3$}
      \end{reponses}
    \end{multicols}
  \end{question}
}
\element{EM1.Quint}{
  \begin{question}{QintCylindreInt}\bareme{1}
    La charge contenue dans une surface de Gauss cylindrique de rayon $r\leq R$ et de hauteur $h$ entourant un cylindre infiniment haut, de rayon $R$ uniformément chargé de densité volumique de charge $\rho$ est
    \begin{multicols}{3}
      \begin{reponses}
        \bonne{$\pi r^2h\rho$}
        \mauvaise{$\pi R^2 h \rho$}
        \mauvaise{$2\pi R^2 h \rho$}
        \mauvaise{$2\pi r^2 h \rho$}
        \mauvaise{$2\pi R h \rho$}
        \mauvaise{$2\pi r h \rho$}
      \end{reponses}
    \end{multicols}
  \end{question}
}
\element{EM1.Quint}{
  \begin{question}{QintBouleInt}\bareme{1}
    La charge contenue dans une surface de Gauss sphérique de rayon $r\geq R$ entourant une boule de rayon $R$ uniformément chargé de densité volumique de charge $\rho$ est
    \begin{multicols}{3}
      \begin{reponses}
        \bonne{$\frac{4}{3}\pi R^3\rho$}
        \mauvaise{$\frac{4}{3}\pi r^3\rho$}
        \mauvaise{$4\pi r^2\rho$}
        \mauvaise{$4\pi R^2\rho$}
        \mauvaise{$4\pi r^3\rho$}
        \mauvaise{$4\pi R^3\rho$}
      \end{reponses}
    \end{multicols}
  \end{question}
}
\element{EM1.Quint}{
  \begin{question}{QintCylindreExt}\bareme{1}
    La charge contenue dans une surface de Gauss cylindrique de rayon $r\geq R$ et de hauteur $h$ entourant un cylindre infiniment haut, de rayon $R$ uniformément chargé de densité volumique de charge $\rho$ est
    \begin{multicols}{3}
      \begin{reponses}
        \mauvaise{$\pi r^2h\rho$}
        \bonne{$\pi R^2 h \rho$}
        \mauvaise{$2\pi R^2 h \rho$}
        \mauvaise{$2\pi r^2 h \rho$}
        \mauvaise{$2\pi R h \rho$}
        \mauvaise{$2\pi r h \rho$}
      \end{reponses}
    \end{multicols}
  \end{question}
}
\setgroupmode{EM1.Quint}{cyclic}

\element{EM1.EPlan}{
  \begin{question}{EPlan}\bareme{1}
    Le champ électrique créé par un plan uniformément chargé, de densité surfacique de charge $\sigma$ s'écrit $\vv{E}=E(z)\vv{e_z}$. On choisi comme surface de Gauss un cylindre de rayon $R$, de hauteur $2z$, symétrique par rapport au plan. Le champ électrique est symétrique par rapport au plan. Le champ électrique s'écrit, en supposant $z>0$
    \begin{multicols}{3}
      \begin{reponses}
        \mauvaise{$\frac{Q_{\text{int}}}{\pi R^2\epsilon_0}\vv{e_z}$}
        \bonne{$\frac{Q_{\text{int}}}{2\pi R^2\epsilon_0}\vv{e_z}$}
        \mauvaise{$\frac{Q_{\text{int}}}{\epsilon_0}\vv{e_z}$}
        \mauvaise{$\frac{Q_{\text{int}}}{\pi R\epsilon_0}\vv{e_z}$}
        \mauvaise{$\frac{Q_{\text{int}}}{2\pi R\epsilon_0}\vv{e_z}$}
        \mauvaise{$\frac{Q_{\text{int}}}{\frac{4}{3}\pi R^2\epsilon_0}\vv{e_z}$}
      \end{reponses}
    \end{multicols}
  \end{question}
}